\section{总结与延伸}

\subsection{核心要点}

\begin{enumerate}
    \item \textbf{AutoGen} 以对话为核心抽象，支持多 Agent 递归组合，适合构建复杂的 Agent 协作系统。
    \item \textbf{LlamaIndex} 从数据质量出发，通过多模态 RAG + Agentic 推理构建生产级知识助手。
    \item 多 Agent 架构相比单 Agent 在任务分解、安全保障和可维护性方面具有显著优势。
    \item 约束型 vs. 通用型 Agent 架构存在可靠性-灵活性权衡。
    \item 生产部署需要考虑微服务化、消息队列和 human-in-the-loop。
\end{enumerate}

\subsection{拓展阅读}

\begin{itemize}
    \item Wu et al., ``AutoGen: Enabling Next-Gen LLM Applications via Multi-Agent Conversation,'' 2023.
    \item Liu, ``LlamaIndex: A Data Framework for LLM Applications,'' 2022--2024.
    \item Lewis et al., ``Retrieval-Augmented Generation for Knowledge-Intensive NLP Tasks,'' NeurIPS 2020.
\end{itemize}

\end{document}
